\documentclass[11pt]{article}

\usepackage[margin=1.5in]{geometry}
\usepackage{lecnotes}
% \usepackage{mpass}
% \lstset{style=verb,language=mpass}
\input{lmacros}

\newcommand{\course}{15-316: Software Foundations of Security \& Privacy}
\newcommand{\lecturer}{Matt Fredrikson}
\newcommand{\lecdate}{March 17, 2026}
\newcommand{\lecnum}{16}
\newcommand{\lectitle}{Authorization Logic}
\newcommand{\courseurl}{https://15316-cmu.github.io/2026/}

\begin{document}

\maketitle

\section{Introduction}

We started the course with an analysis of \emph{safety properties} of
computations, and how to account for them statically (via verification conditions in
dynamic logic) or dynamically (via sandboxing).  Then we moved on to the more
complex \emph{information flow properties}.  An attacker may discover secrets
via particular programs or inputs in a variety of ways, with side-channel
attacks (for example, timing attacks) being the most sophisticated.
Countermeasures are mostly via information-flow type systems.

With this lecture we are starting a new section of the course, considering
\emph{authentication} (you are who you say you are) and \emph{authorization}
(you are allowed to perform the actions you are trying to perform).
Authentication and authorization are pervasive in today's computing environment,
from shared file systems like AFS and cloud services like Github, to shopping
and banking services.  Today's and the next lecture will focus on
\emph{authorization}, followed in later lectures by authentication.  In many
cases, authorization is just embedded in code.  Once authorization becomes
complex, this can easily be compromised by bugs since flow of the authorization
checks through a program can be difficult to audit.  It can also be difficult to
even understand what the authorization policy actually is, which then means it
is hard to compare it against the code.

In this lecture we take a very general approach, expressing policy in an
\emph{authorization logic}.  Like in other applications we have seen, the logic
is the interface between policy and implementation.  It expresses, at a high
level of abstraction, who is allowed to do what in a complex system.  On one
side, this serves as a specification one can reason about rigorously, divorced
from an implementation.  On the other side, we can use it directly to enforce
authorization policies in an implementation by using formal proofs of
authorization.  \citet{Abadi03lics} gives a general overview of authorization
logics.  The architecture where formal proofs are used was pioneered by
\citet{Bauer03phd} under the name of \emph{proof-carrying authorization} (PCA).

Among the direct applications of PCA are the Grey systems~\citep{Bauer05isc} for
access to offices in Cylab at CMU, and a \emph{proof-carrying file system} that
particularly explored issues of efficiency~\citep{Garg09phd,Garg10oakland}.  We
ignore here a number of practically relevant issues, such as revocation and
temporal aspects of authorization.  Once added~\citep{DeYoung07tr},
authorization logic is expressive enough, for example, to express the rules
governing access to classified information in the American intelligence
community~\citep{Garg09trpolicy}.

\section{Affirmations}

Early on this class, we introduced and reasoning in Boolean logic, with
connectives such as conjunction, disjunction, implication, etc.  Then we
introduced so-called \emph{modal operators}, $[\alpha]Q$,
$\langle\alpha\rangle Q$ that speak about programs, and $\Box P$ that guarantees
validity of $P$ (that is, truth in all possible states).

In order to express access control policies \emph{logically}, we need a new kind
of modal operator that expresses an affirmation, $A \says P$ (\emph{principal
  $A$ says proposition $P$}).  Principals $A$, $B$, $C$, etc.\ can stand for
user ids like $\mi{admin}$, $\mi{mfredrik}$, or $\mi{saranyav}$. Propositions include
atomic propositions, conjunction, implication, quantifiers, and other
connectives as we need them.

As an example, we use the Grey~\citep{Bauer05isc} system that is used to control
access to offices in Cylab at CMU\@.  There is an administrator (principal
$\mi{admin}$) that sets policies, and individual professors and students
identified by their Andrew id.  There are also resources, like offices and
conference rooms.  Among the atomic propositions are the following:
\begin{itemize}
\item $\ms{mayOpen}(A, R)$.  Principal $A$ may open room $R$.
\item $\ms{owns}(A, R)$.  Principal $A$ owns office $R$.
\item $\ms{studentOf}(B, A)$.  Principal $B$ is a student of $A$.
\end{itemize}
Here are some examples.  The administrator may say that $\mi{mfredrik}$ owns
$\mi{cic2126}$.
\begin{tabbing}
  \quad $\mi{admin} \says \ms{owns}(\mi{mfredrik}, \mi{cic2126})$
\end{tabbing}
This is a basic affirmation.  A general rule could state that the owner of a
room may open it.
\begin{tabbing}
  \quad $\mi{admin} \says (\forall A.\, \forall R.\, \ms{owns}(A, R) \arrow \ms{mayOpen}(A, R))$
\end{tabbing}
An even more complex policy component would be that any student of the owner of
an office, may also open the office.  The twist here is that the
$\ms{studentOf}$ relationship should be affirmed by the owner, rather than the
administrator.
\begin{tabbing}
  \quad $\mi{admin} \says (\forall A.\, \forall B.\, \forall R.\, \ms{owns}(A, R)
  \land A \says \ms{studentOf}(B, A) \arrow \ms{mayOpen}(B, R))$
\end{tabbing}
The scope of $A \says$ here is intended to be only the $\ms{studentOf}$
proposition.  When the scope is larger, we enclose it in parentheses.  We can
combine this policy with the affirmation by $\mi{mfredrik}$ that $\mi{saranyav}$ is
his student:
\begin{tabbing}
  \quad $\mi{mfredrik} \says \ms{studentOf}(\mi{saranyav}, \mi{mfredrik})$
\end{tabbing}
Under this policy $\mi{saranyav}$ should be able to access $\mi{cic2126}$.  We can
formulate the question whether this is allowed as a sequent in authorization
logic, with the policy as antecedents (assumptions) and the query as the
succedent.  In this particular situation, we would try to prove the sequent
\begin{tabbing}
  \quad $(1) : \mi{admin} \says (\forall A.\, \forall R.\, \ms{owns}(A, R) \arrow \ms{mayOpen}(A, R))$, \\
  \quad $(2) : \mi{admin} \says (\forall A.\, \forall B.\, \forall R.\, \ms{owns}(A, R)
  \land A \says \ms{studentOf}(B, A) \arrow \ms{mayOpen}(B, R))$, \\[1ex]
  \quad $(3) : \mi{admin} \says \ms{owns}(\mi{mfredrik}, \mi{cic2126})$, \\
  \quad $(4) : \mi{mfredrik} \says \ms{studentOf}(\mi{saranyav}, \mi{mfredrik})$ \\[1ex]
  \quad $\vdash$ \\[1ex]
  \quad $\mi{admin} \says \ms{mayOpen}(\mi{saranyav}, \mi{cic2126})$
\end{tabbing}
Here, we have labeled the assumptions so we can reference them in the
proof.  Reasoning entirely intuitively (to be formalized later in the
lecture) we can deduce:
\begin{tabbing}
  \quad $(5) : \mi{admin} \says \ms{mayOpen}(\mi{mfredrik}, \mi{cic2126})$ \` from $(1)$ and $(3)$ \\
  \quad $(6) : \mi{admin} \says \ms{mayOpen}(\mi{saranyav}, \mi{cic2126})$ \` from $(2)$, $(3)$, and $(4)$
\end{tabbing}
The last line $(6)$ is exactly what we are trying to prove.  So we have
confirmed that $\mi{saranyav}$ may open my office, according to the policy.

\section{Constructive Logic}

The Boolean logic of the earlier lectures (including propositional logic and
dynamic logic) are based on a semantics where every formula has one of two truth
values: $\top$ or $\bot$.  This is not a good match for authorization logic, as
remarked by \citet{Abadi03lics}.  The first issue that we may be ``agnostic''
about a proposition.  Any proposition that is affirmed should be seen as
\emph{extending} what we can deduce, but not be in conflict with what we know so
far.  Another example is given by
\[
  A \says P \arrow (P \lor A \says Q)
\]
We can read this: if $A$ affirms $P$ then either $P$ must be true, or $A$
affirms any proposition (including $\bot$).  If we worked with just two truth
values, why would this be valid?  Let's assume $A \says P$.  Now we distinguish
two cases for $P$.  If $P$ is true, then $P \lor A \says Q$.  If $P$ is false,
then $\lnot P$ is true.  $A$ would affirm any true proposition, so
$A \says \lnot P$.  But if $A \says P$ and $A \says \lnot P$, $A$ is mired in a
state of internal contradiction and would have to affirm anything.  Not good.

In order to avoid these kind of paradoxes, we use an \emph{intutionistic logic}
as the basis for reasoning.  Intuitionistic logic does not allow us to reason
with the law of excluded middle, or to prove $P$ by assuming $\lnot P$ and
deriving a contradiction (using an indirect proof).  Instead, we want the reason
access to a resource is granted to be as clear and direct as possible, which is
what intuitionistic logic provides.  Actually, we go a step further and
rule out negation $\lnot P$ and falsehood $\bot$ entirely, because it is easy to
make intuitively meaningful statements that are wrong.  For example, principal
$\mi{mfredrik}$ does \textbf{not} want $\mi{saranyav}$ to access his office.  Stating
$\mi{mfredrik} \says \lnot \ms{mayOpen}(\mi{saranyav}, \mi{cic2126})$ turns out to be
the wrong way to say this.  Because if the rest of the policy says that
$\ms{mayOpen}(\mi{saranyav}, \mi{cic2126})$ then suddenly $\mi{mfredrik}$ affirms
\emph{everything}, including that every principal in the system may access his
office---clearly not the desired effect.

This change in perspective, from two-valued Boolean logic (also called
\emph{classical logic}) to a richer intuitionstic logic has two consequences,
one semantic and one syntactic.  Semantically, we need to generalize to a
so-called \emph{Kripke semantics} where we consider multiple \emph{worlds} in
which different propositions may be true.  This is a good match for a logic of
authorization since at the very least each principal defines a world, and
different principals will affirm different propositions.  For such a semantics,
see, for example, \citet{Garg08imla}.  Syntactically, it is surprisingly simple: we
restrict the succedent of a sequent to be exactly one formula.  We therefore
write $\Gamma \vdash \delta$.  That this actually works was one of
Gentzen's~\citeyearpar{Gentzen35} profound insights.

We actually make the lack of a mathematical semantics a philosophical principle.
We think of the meanings of formulas as given by their possible derivations in
the sequent calculus.  In other words, the right and left rules of the sequent
calculus themselves define the meaning of each logical constant, connective, and
modality.  In the context of authorization, this can be justified since it is
ultimately a formal proof that is used to claim and check authorization---we
don't appeal to an external semantics.  For a fuller development of this
viewpoint, see, for example, \citet{Dummett91} and \citet{MartinLof83}.

As we will in particular see in the next lecture, the properties of the logic
change in fundamental ways when we restrict the succedents to be just a single
formula.  The rules we get otherwise just reflect the earlier ones.
\begin{rules}
  \infer[\ms{id}]
  {\Gamma, P \vdash P}
  {}
  \\[1em]
  \infer[{\arrow}R]
  {\Gamma \vdash P \arrow Q}
  {\Gamma, P \vdash Q}
  \hspace{3em}
  \infer[{\arrow}L]
  {\Gamma, P \arrow Q \vdash \delta}
  {\Gamma \vdash P
    & \Gamma, Q \vdash \delta}
  \\[1em]
  \infer[{\land}R]
  {\Gamma \vdash P \land Q}
  {\Gamma \vdash P
    & \Gamma \vdash Q}
  \hspace{3em}
  \infer[{\land}L]
  {\Gamma, P \land Q \vdash \delta}
  {\Gamma, P, Q \vdash \delta}
  \\[1em]
  \infer[{\forall}R^y]
  {\Gamma \vdash \forall x.\, P(x)}
  {\Gamma \vdash P(y) & \mbox {$y \not\in (\Gamma, \forall x.\, P(x))$}}
  \hspace{3em}
  \infer[{\forall}L]
  {\Gamma, \forall x.\, P(x) \vdash \delta}
  {\Gamma, P(c) \vdash \delta}
  \\[1em]
  \infer[{\lor}R_1]
  {\Gamma \vdash P \lor Q}
  {\Gamma \vdash P}
  \hspace{1em}
  \infer[{\lor}R_2]
  {\Gamma \vdash P \lor Q}
  {\Gamma \vdash Q}
  \hspace{3em}
  \infer[{\lor}L]
  {\Gamma, P \lor Q \vdash \delta}
  {\Gamma, P \vdash \delta
    & \Gamma, Q \vdash \delta}
\end{rules}
For the quantifiers, recall our convention of writing $P(x)$ and then $P(c)$ for
the result of substituting $c$ for $x$ in $P(x)$.  Quantification was previously
over integers, but here we think of principals, rooms, etc.  $c$ could also be a
variable $y$ that was introduced by a ${\forall}R$ rule.

The most immediate effect of reasoning intuitionistically is perhaps on
disjunction.  We can no longer prove, for example $p \lor (p \arrow q)$ because
we have to decide between one of the disjuncts instead of carrying both as
succedents.

In these rules there is a main formula among the antecedents to which a left
rule is applied.  This formula may be needed again so it can be kept among the
antecedents if so desired.  This is particularly useful for ${\forall}L$.  To
anticipate a bit the next lecture, it is no longer the case that all rules are
invertible, the way it is in Boolean propositional logic.  So proof search is
significantly more difficult than in Boolean logic, and not just because of the
quantifiers.

\section{Affirmations}
\label{sec:affirmations}

We haven't yet discussed $A \says P$, the central modality of authorization
logic.  The key insight is that when we try to prove $\Gamma \vdash A \says P$
we need to proceed with our reasoning from the perspective of principal $A$.
But what does this mean?  First, principal $A$ should be willing to affirm any
proposition that is \emph{true}.  Second, anything that $A$ says is also
available to reason with.  On the other hand, if $B \says Q$ is in $\Gamma$ for
some $B \neq A$, the proposition $Q$ is not available to $A$: this is something
that $B$ affirms, but not necessarily $A$.

Formalizing this reasoning is not entirely straightforward, and there are
different approaches.  The one we choose is due to \citet{Garg06csfw}, primarily
due to it simplicity.  We distinguish a succedent expressing that a proposition
is true (written $P\; \ms{true}$) and another that $A$ affirms a proposition $P$
(written $A \aff P$).  We often omit the ``$\ms{true}$'' notation, but always
write $A \aff P$ to use an affirmation.  Formally:
\[
  \begin{array}{llcl}
    \mbox{Succedent} & \delta & ::= & P\; \ms{true} \mid A \aff P
  \end{array}
\]
The first rule says that in order to prove that $A \says P$ is true,
we need to prove that $A$ affirms $P$.
\begin{rules}
  \infer[{\mb{says}}R]
  {\Gamma \vdash (A \says P)\; \ms{true}}
  {\Gamma \vdash A \aff P}
\end{rules}
This may seem redundant, but it exposes the principal $A$ and the proposition
$P$.  It is similar to the rule ${\arrow}R$, where the implication $P \arrow Q$
is turned into $P \vdash Q$, opening up $P$ and $Q$ to further inferences.

The next rule expresses that if $P$ is true, then $A$ is willing to affirm that.
\begin{rules}
  \infer[\mb{aff}]
  {\Gamma \vdash A \aff P}
  {\Gamma \vdash P\; \ms{true}}
\end{rules}
This rule is like ``peeling an onion'', moving entirely now into $A$'s head,
reasoning from their perspective.

The left rule for $A \says P$ jumps directly to $P$, both of which are
propositions and therefore can appear among the antecedents.  But we can
make this transition only if we are currently reasoning from $A$'s perspective,
that is, if the succedent is $A \aff Q$ for some $Q$.
\begin{rules}
  \infer[{\mb{says}}L]
  {\Gamma, A \says P \vdash A \aff Q}
  {\Gamma, P \vdash A \aff Q}
\end{rules}

\section{Some Axioms}

Before we go back to our motivating example, we can analyze some of the
properties of affirmations.
\begin{tabbing}
  \qquad $\vdash P \arrow A \says P$
\end{tabbing}
As we have said, any principal (here $A$) is willing to affirm any true
proposition (here $P$), so we should be able to prove this.  As usual,
we build the proof bottom-up; we show here only the result.
\[
  \infer[{\arrow}R]
  {\cdot \vdash P \arrow A \says P}
  {\infer[{\says}R]
    {P \vdash A \says P}
    {\infer[{\aff}]
      {P \vdash A \aff P}
      {\infer[\ms{id}]{P \vdash P}{}}}}
\]
The implication in the other direction should not be valid: just because
$A$ says $P$ that doesn't mean $P$ is actually true.  We show that it is not
derivable.
\begin{tabbing}
  \qquad $\not\vdash (A \says p) \arrow p$
\end{tabbing}
\[
  \infer[{\arrow}R]
  {\cdot \vdash (A \says p) \arrow p}
  {\deduce[\ms{XXX}]{A \says p \vdash p}{}}
\]
As a first step, only ${\arrow}R$ is applicable; for the second no rule is.

We can also ``distribute'' $A \says \null$ over an implication.  This captures
that if $A$ affirms $P \arrow Q$ and $P$, then it also affirms $Q$.
\begin{tabbing}
  \qquad $\vdash A \says (P \arrow Q) \arrow (A \says P \arrow A \says Q)$
\end{tabbing}
\[
  \infer[{\arrow}R \times 2]
  {\vdash A \says (P \arrow Q) \arrow (A \says P \arrow A \says Q)}
  {\infer[{\says}R]
    {A \says (P \arrow Q), A \says P \vdash A \says Q}
    {\infer[{\says}L \times 2]
      {A \says (P \arrow Q), A \says P \vdash A \aff Q}
      {\infer[\mb{aff}]
        {P \arrow Q, P \vdash A \aff Q}
        {\infer[{\arrow}L]
          {P \arrow Q, P \vdash Q}
          {\infer[\ms{id}]{P \vdash P}{}
            & \infer[\ms{id}]{Q \vdash Q}{}}}}}}
\]
Iterating $A \says \null$ twice is the same as just affirming just once.
\begin{tabbing}
  \qquad $\vdash A \says (A \says P) \arrow A \says P$
\end{tabbing}
\[
  \infer[{\arrow}R]
  {\vdash A \says (A \says P) \arrow A \says P}
  {\infer[{\says}R]
    {A \says (A \says P) \vdash A \says P}
    {\infer[{\says}L]
      {A \says (A \says P) \vdash A \aff P}
      {\infer[{\says}L]
        {A \says P \vdash A \aff P}
        {\infer[\aff]
          {P \vdash A \aff P}
          {\infer[\ms{id}]{P \vdash P}{}}}}}}
\]
The other direction of this implication is an instance of the first axiom.  On
the other hand, for different principals $A$ and $B$ we cannot prove
\[
  \qquad \not\vdash A \says p \arrow B \says p
\]
These axioms (together with those for intuitionistic logic) are complete for
implication and affirmation and identify this as a generalization of \emph{lax
  logic} \citep{Fairtlough97ic}.  Instead of a single modality $\bigcirc P$ we
have a whole family of such modalities, indexed by principals.  From the
perspective of functional programming, each modality ``$A \says {-}$'' is a strong
monad \citep{Moggi89,Wadler92popl}.  Here, it relativizes reasoning to each
principal; in functional programming monads can isolate the pure part of the
language from effects.  The connection to modal logics is further explored by
\citet{Pfenning01mscs}.  A summary of the axioms is in \autoref{fig:axioms}.

\begin{figure}[ht]
  \begin{tabbing}
  \qquad $\vdash P \arrow A \says P$ \\
  \qquad $\vdash A \says (P \arrow Q) \arrow (A \says P \arrow A \says Q)$ \\
  \qquad $\vdash A \says (A \says P) \arrow A \says P$ \\[1em]
  \qquad $\not\vdash (A \says p) \arrow p$
  \end{tabbing}
  \caption{Axioms for Authorization Logic}
  \label{fig:axioms}
\end{figure}

\section{An Example of Authorization}

We return to our motivating example, where we have labeled each of the
antecedents as before.
\begin{tabbing}
  \quad $(1) : \mi{admin} \says (\forall A.\, \forall R.\, \ms{owns}(A, R) \arrow \ms{mayOpen}(A, R))$, \\
  \quad $(2) : \mi{admin} \says (\forall A.\, \forall B.\, \forall R.\, \ms{owns}(A, R)
  \land A \says \ms{studentOf}(B, A) \arrow \ms{mayOpen}(B, R))$, \\[1ex]
  \quad $(3) : \mi{admin} \says \ms{owns}(\mi{mfredrik}, \mi{cic2126})$, \\
  \quad $(4) : \mi{mfredrik} \says \ms{studentOf}(\mi{saranyav}, \mi{mfredrik})$ \\[1ex]
  \quad $\vdash$ \\[1ex]
  \quad ${\color{blue}\mi{admin} \says\null} \ms{mayOpen}(\mi{saranyav}, \mi{cic2126})$
\end{tabbing}
We highlighted in blue the focus of the next inference.  Using $\mb{says}R$, we
reduce this in one step to proving the sequent
\begin{tabbing}
  \quad $(1) : {\color{blue}\mi{admin} \says\null} (\forall A.\, \forall R.\, \ms{owns}(A, R) \arrow \ms{mayOpen}(A, R))$, \\
  \quad $(2) : {\color{blue}\mi{admin} \says\null} (\forall A.\, \forall B.\, \forall R.\, \ms{owns}(A, R)
  \land A \says \ms{studentOf}(B, A) \arrow \ms{mayOpen}(B, R))$, \\[1ex]
  \quad $(3) : {\color{blue}\mi{admin} \says\null} \ms{owns}(\mi{mfredrik}, \mi{cic2126})$, \\
  \quad $(4) : \mi{mfredrik} \says \ms{studentOf}(\mi{saranyav}, \mi{mfredrik})$ \\[1ex]
  \quad $\vdash$ \\[1ex]
  \quad $\mi{admin} \aff \ms{mayOpen}(\mi{saranyav}, \mi{cic2126})$
\end{tabbing}
This unlocks the affirmations by $\mi{admin}$ among the antecedents, so
applying $\mb{says}L$ three times we get
\begin{tabbing}
  \quad $(1)' : \forall A.\, \forall R.\, \ms{owns}(A, R) \arrow \ms{mayOpen}(A, R)$, \\
  \quad $(2)' : {\color{blue}\forall A.\, \forall B.\, \forall R.}\, \ms{owns}(A, R)
  \land A \says \ms{studentOf}(B, A) \arrow \ms{mayOpen}(B, R)$, \\[1ex]
  \quad $(3)' : \ms{owns}(\mi{mfredrik}, \mi{cic2126})$, \\
  \quad $(4) : \mi{mfredrik} \says \ms{studentOf}(\mi{saranyav}, \mi{mfredrik})$ \\[1ex]
  \quad $\vdash$ \\[1ex]
  \quad $\mi{admin} \aff \ms{mayOpen}(\mi{saranyav}, \mi{cic2126})$
\end{tabbing}
Note that $\mi{mfredrik}$'s affirmation cannot be unlocked, since
$\mi{admin} \neq \mi{mfredrik}$.  Now we can apply ${\forall}L$ three times, instantiating
$A$, $B$, and $R$ with $\mi{mfredrik}$, $\mi{saranyav}$, and $\mi{cic2126}$, respectively.
\begin{tabbing}
  \quad $(1)' : \forall A.\, \forall R.\, \ms{owns}(A, R) \arrow \ms{mayOpen}(A, R)$, \\
  \quad $(2)'' : {\color{blue}\ms{owns}(\mi{mfredrik}, \mi{cic2126})
  \land \mi{mfredrik} \says \ms{studentOf}(\mi{saranyav}, \mi{mfredrik})} \arrow \ms{mayOpen}(\mi{saranyav}, \mi{cic2126})$, \\[1ex]
  \quad $(3)' : {\color{blue}\ms{owns}(\mi{mfredrik}, \mi{cic2126})}$, \\
  \quad $(4) : {\color{blue}\mi{mfredrik} \says \ms{studentOf}(\mi{saranyav}, \mi{mfredrik})}$ \\[1ex]
  \quad $\vdash$ \\[1ex]
  \quad $\mi{admin} \aff \ms{mayOpen}(\mi{saranyav}, \mi{cic2126})$
\end{tabbing}
At this point we can apply ${\arrow}L$ to the antecedent $(2)''$, proving
the conjunction by $(3)'$ and $(4)$ and obtaining the new line $(5)$.
\begin{tabbing}
  \quad $(1)' : \forall A.\, \forall R.\, \ms{owns}(A, R) \arrow \ms{mayOpen}(A, R)$, \\
  \quad $(2)'' : \ms{owns}(\mi{mfredrik}, \mi{cic2126})
  \land \mi{mfredrik} \says \ms{studentOf}(\mi{saranyav}, \mi{mfredrik}) \arrow \ms{mayOpen}(\mi{saranyav}, \mi{cic2126})$, \\[1ex]
  \quad $(3)' : \ms{owns}(\mi{mfredrik}, \mi{cic2126})$, \\
  \quad $(4) : \mi{mfredrik} \says \ms{studentOf}(\mi{saranyav}, \mi{mfredrik})$ \\
  \quad $(5) : {\color{blue}\ms{mayOpen}(\mi{saranyav}, \mi{cic2126})}$ \\[1ex]
  \quad $\vdash$ \\[1ex]
  \quad ${\color{blue}\mi{admin} \aff \ms{mayOpen}(\mi{saranyav}, \mi{cic2126})}$
\end{tabbing}
Now we can apply the rule of affirmation, followed by the identity to complete
the proof.

\section{Summary}

\[
  \begin{array}{llcl}
    \mbox{Principals} & A, B, C \\
    \mbox{Formulas} & P, Q & ::= & p \mid P \land Q \mid P \arrow Q \mid P \lor Q \mid \forall x.\, P(x) \mid A \says P \\
    \mbox{Succedents} & \delta & ::= & P\; \ms{true} \mid A \aff P
  \end{array}
\]

\begin{figure}[ht]
  \begin{rules}
  \infer[\ms{id}]
  {\Gamma, P \vdash P\; \ms{true}}
  {}
  \\[1em]\hline\\[1ex]
  \infer[{\arrow}R]
  {\Gamma \vdash P \arrow Q\; \ms{true}}
  {\Gamma, P \vdash Q\; \ms{true}}
  \hspace{3em}
  \infer[{\arrow}L]
  {\Gamma, P \arrow Q \vdash \delta}
  {\Gamma \vdash P
    & \Gamma, Q \vdash \delta}
  \\[1em]
  \infer[{\land}R]
  {\Gamma \vdash P \land Q\; \ms{true}}
  {\Gamma \vdash P\; \ms{true}
    & \Gamma \vdash Q\; \ms{true}}
  \hspace{3em}
  \infer[{\land}L]
  {\Gamma, P \land Q \vdash \delta}
  {\Gamma, P, Q \vdash \delta}
  \\[1em]
  \infer[{\forall}R^y]
  {\Gamma \vdash \forall x.\, P(x)\; \ms{true}}
  {\Gamma \vdash P(y)\; \ms{true} & \mbox {$y \not\in (\Gamma, \forall x.\, P(x))$}}
  \hspace{3em}
  \infer[{\forall}L]
  {\Gamma, \forall x.\, P(x) \vdash \delta}
  {\Gamma, P(c) \vdash \delta}
  \\[1em]
  \infer[{\lor}R_1]
  {\Gamma \vdash P \lor Q\; \ms{true}}
  {\Gamma \vdash P\; \ms{true}}
  \hspace{1em}
  \infer[{\lor}R_2]
  {\Gamma \vdash P \lor Q\; \ms{true}}
  {\Gamma \vdash Q\; \ms{true}}
  \hspace{3em}
  \infer[{\lor}L]
  {\Gamma, P \lor Q \vdash \delta}
  {\Gamma, P \vdash \delta
    & \Gamma, Q \vdash \delta}
  \\[1em]\hline\\[1ex]
  \infer[{\mb{says}}R]
  {\Gamma \vdash (A \says P)\; \ms{true}}
  {\Gamma \vdash A \aff P}
  \hspace{3em}
  \infer[{\mb{says}}L]
  {\Gamma, A \says P \vdash A \aff Q}
  {\Gamma, P \vdash A \aff Q}
  \\[1em]
  \infer[\mb{aff}]
  {\Gamma \vdash A \aff P}
  {\Gamma \vdash P\; \ms{true}}
  \end{rules}
  \caption{Affirmation Logic}
\end{figure}


\bibliographystyle{plainnat}
\bibliography{bibliography}

\end{document}
